% 本文件由 LaTeX 排版 Agent 根据有效 translation.json 自动生成。
% author=Sulpitius Severus (c. 363-c. 420); work_id=3501; title=论圣玛尔定生平

\chapter{论\Person{圣玛尔定}生平}

% structure: p0001 source=h1 target=omitted reason=page-title-already-rendered-by-parent

% structure: p0002 source=h2 target=section reason=relative-heading-rank; base=h2

\section{致\Person{德西德留}序言}

\Person{塞维鲁}致他最亲爱的兄弟\Person{德西德留}，祝安。志同道合的兄弟，我曾决定将我所写的关于\Person{圣玛尔定}生平的小册子秘藏于己室，不示外人。我这样做，是因为我并无出众的才华，且畏惧世人的批评，恐怕（如我所料）我那略显粗拙的文风会使读者不悦，而我因过于大胆地选取了一个本当留给真正雄辩的作家去写的题材，会被人视为理应受到普遍的责备。但我无法拒绝你一再提出的请求。因着对你的爱，有什么事是我不会应允的呢，即使以牺牲我的谦逊为代价？然而，我将此作交付给你，是建立在确知你不会将其透露给任何人的基础上，因为你已向我作出承诺。尽管如此，我仍然担心你会成为它公诸于世的中介；我深知一旦发表，便永不可收回。若此事发生，且你得知它被某些人阅读，我希望你能善意地请求读者留意所叙述的事实，而非表达这些事实的语言。你要恳请他们，若文风不合耳，勿要见怪，因为天主的国不在于雄辩，而在于信德。也请他们记住，救恩并非由演说家，而是由渔夫传报给世界，尽管天主若有益处，无疑也可采取另一途径。就我而言，当我最初决意撰写下文时，因我认为如此伟大人物的卓越事迹若被隐没是件可耻的事，我便下定决心不为语言的语病而感到羞耻。我这样做，是因为我从未在这方面获得任何高深的知识；或者，即使我从前曾对此类学问稍有涉猎，但因长期荒废，也已尽数遗忘。然而，为了将来不必再采取这种令人厌烦的自辩方式，最好的办法是：若你赞成，此书隐去作者姓名出版。为此，请将书册封面上的标题擦去，使扉页保持沉默；（这已足够）让本书只宣告其主题，而不透露作者是谁。\par

% structure: p0004 source=h2 target=section reason=relative-heading-rank; base=h2

\section{第一章}

撰写\Person{圣玛尔定}生平的理由\par

大多数人徒然追求世俗的荣耀，并以为由此获得了自己名字的纪念，即：以笔墨来美化名人的生平。这种做法，虽然并未为他们赢得持久的声誉，却无疑带来了他们所怀希望的一些实现。这样做，既徒劳地保存了他们自己的记忆，又因将伟人的榜样呈献给世人，在读者心中激起了不小的竞效之情。然而，尽管如此，他们的劳苦丝毫未能关乎我们所期盼的永福和永生的生命。因为仅仅注定随世界一同消亡的荣耀，对那些只写世俗事物的人又有什么益处呢？后世从阅读\Person{赫克托耳}作为战士或\Person{苏格拉底}作为哲学阐释者的事迹中又得到什么好处呢？这些事毫无益处，因为不仅效法所提及的人物是愚蠢的，而且不以最大的严厉抨击他们，简直是疯狂。事实上，那些仅凭现世行为来评估人生的人，已将他们的希望寄托于虚妄，将他们的灵魂送入了坟墓。他们仅令自己存留于必死者的记忆里，而人的本分却是追求永生，而非永久的纪念，且不是通过写作、战斗或哲学思考，而是通过度虔诚、圣洁、宗教的生活。人类的这种错误行为，因被载入文学而盛行，以致有众多人效法那被称颂的虚妄哲学或愚昧的卓越。因此，我认为若能写出一位至圣者的生平，作为将来他人的榜样，那将是一件值得付出必要辛劳的成就；事实上，这将唤醒读者追求真知、属天的争战和神圣的德行。这样做，我们也顾念自己的益处，期望从天主那里得到的，不是人间的虚名，而是永恒的赏报。因为，虽然我们自己并未如此生活，以至于能成为他人的榜样，但我们仍力求使那位值得效法的人不被隐没。因此，我将着手撰写\Person{圣玛尔定}的生平，叙述他任主教前以及任主教时所行的事。同时，我不敢奢望能述尽他的一切所是或所行。那些唯独他自己知晓的卓越品行，我们全然不知，因为他并不寻求世人的荣誉，反而尽量隐藏自己的美德。即使那些为我们所知的事迹，我们也省略了许多，因为我们认为只记录那些较突出、卓越的事迹便已足够。同时，为了读者的利益，我必须注意，不要因列举过多的事例而使他们生厌。但我恳请将要阅读下文的人，完全相信所叙述的事，并相信我没有写任何我无确切知识和证据的事。事实上，我宁愿保持沉默，也不愿叙述虚假的事。\par

% structure: p0007 source=h2 target=section reason=relative-heading-rank; base=h2

\section{第二章}

\Person{圣玛尔定}的兵役\par

玛尔定出生于潘诺尼亚的\Place{萨巴里亚}，但在意大利的\Place{提契努姆}长大。按世俗的标准，他的父母并非微贱之辈，但他们是外教人。他的父亲最初是一名普通士兵，后来成为护民官。玛尔定本人年轻时追随军旅生涯，被编入帝国卫队，先后在君士坦丁王和尤利安凯撒麾下服役。然而，这并非出于他自己的意愿，因为几乎从幼年起，这位杰出少年的圣洁幼年便渴慕事奉天主。当他十岁时，不顾父母的意愿，自行前往教会，恳求成为慕道者。不久之后，他奇妙地完全献身于事奉天主，十二岁时便渴望度隐修者的生活；他本会以必要的誓愿来实现这一愿望，若非他尚过于年轻而受阻。然而，他的心始终专注于隐修院或教会的事务，早在少年时代就已默想他后来作为基督的公开仆人所成就的事。但当国家统治当局颁布法令，要求老兵之子应征入伍时，他因父亲（嫉妒他的圣善行为）的告发而被捕并被锁链捆绑，十五岁时被迫宣誓参军。那时，他仅满足于一名仆从侍候，甚至常常与仆从互换角色，尽管他是真正的主人，却表现得更像下属，以至于时常亲手为仆从脱靴并擦净；他们一同用餐，但真正的主人通常扮演仆人的角色。在领受圣洗前的近三年中，他从事军旅，但完全远离了那类人常陷入的恶习。他对同袍极其和善，怀有奇妙的友爱；而他的忍耐与谦逊超越了人性所能及。无需赞扬他所表现出的克己：它如此伟大，以至于在那时就已不被视为士兵，而更像一位隐修士。凭借所有这些品质，他赢得了全体战友的爱戴，他们既钦佩他又奇妙地爱他。虽然尚未在基督内成为新受造物，他却以善行扮演了慕道者的角色。例如，他帮助困苦者，资助不幸者，支持贫穷者，赤身者，除每日必需外，他未从军饷中为自己保留任何东西。即使在那时，他远非漫不经心地听福音，而是如此遵守其诫命，以至于不为明天忧虑。\par

% structure: p0010 source=h2 target=section reason=relative-heading-rank; base=h2

\section{第三章}

基督显现给圣玛尔定\par

因此，在一个特定的时期，当他除武器和简单的军装外一无所有时，适逢隆冬，一个比往常更为严寒的冬天，以致极度的寒冷使许多人丧命。他恰好在亚眠城门口遇到一个赤身露体的穷人。那人恳求路人怜悯他，但众人都没有留意这个可怜人。这时，这位充满天主的玛尔定认出，那个无人怜悯的人，在这一点上是留给了他的。但他该怎么办？他只有身上穿着的外衣，因为他已经把其余衣物为类似目的施舍了。于是，他拿起佩带的剑，将外衣分成两半，一半给了穷人，自己又穿上了另一半。对此，一些旁观者嘲笑他，因为他现在成了一个难看的对象，只穿着半件衣服。然而，许多有更健全判断力的人深深叹息，因为他们自己没有做过类似的事。他们尤其感到遗憾，因为他们拥有的比玛尔定多，本可以施舍穷人而不至于赤身。在随后的夜晚，当玛尔定入睡时，他看见基督穿着他给穷人的那半件外衣。他极其仔细地凝视着主，并被告知要承认他施舍的袍子是属于他的。不久，他听见耶稣以清晰的声音对周围的众多天使说：\Quote{玛尔定，他仍只是个慕道者，用这件外衣披在了我身上\BibleRef{玛 25:40}。}主确实记念自己的话（他曾在世上说过：\Quote{你们对这些最小兄弟中的一个所做的，就是对我做的。}），宣称他自己就是在那穷人身上被穿上了衣服；为了证实他对如此善行的见证，他屈尊让玛尔定看见他自己正穿着那穷人接受的衣服。在这神视之后，这位圣人并未被人间的荣耀所骄傲，反而承认天主在此事中所行的美善，并因已年满二十岁，急忙领受了圣洗。然而，他并未立即退出军旅，因为他的护民官恳求他留下，而玛尔定允许他成为自己亲密的朋友和同帐篷的伙伴。护民官许诺，在他任期届满后，他也要退出世俗。玛尔定因期待此事而被挽留，在领洗后近两年中，虽名义上仍扮演士兵的角色。\par

% structure: p0013 source=h2 target=section reason=relative-heading-rank; base=h2

\section{第四章}

玛尔定退出军旅服务\par

与此同时，野蛮人涌入\Place{高卢}的两部分，\Person{尤利安凯撒}在\Place{沃季奥内斯城}集结军队，开始向士兵分发赏金。按照惯例，他们被依次叫上前去，直到轮到\Person{马丁}。这时，马丁认为这是请求退伍的适当时机——因为他觉得，如果不再继续服役，就不应接受赏金——于是对凯撒说：\Quote{迄今为止，我作为士兵为\Emph{你}服役；现在请允许我成为天主的士兵：让那将要服侍你的人接受你的赏金吧；我是基督的士兵：我不可以战斗。}那暴君听到这些话勃然大怒，宣称马丁是因为害怕第二天即将到来的战斗，而非出于宗教情感才退出军役。但马丁满怀勇气，甚至因眼前的危险而更加坚定，高喊道：\Quote{如果我的行为被归因于怯懦而非信德，那我明天将不带武器站在战线前，因主耶稣之名，受十字架标记的保护，而非盾牌或头盔，我将安然穿过敌阵。}因此，他被下令关回监狱，决心以无武装面对野蛮人来证明自己所言非虚。但第二天，敌人派遣使节商议和平，并交出他们自己和所有财产。在这种情况下，谁能怀疑这胜利是出于那位圣人呢？他被赐予恩宠，不必无武装上阵。尽管仁慈的主本可即使在敌人的刀剑和标枪中保护他自己的士兵，但为使他的福眼不因目睹他人死亡而痛苦，他消除了战斗的必要。因为基督不需要为自己的士兵赢得任何别的胜利，只需敌人不流血地被制伏，无人丧命。\par

% structure: p0016 source=h2 target=section reason=relative-heading-rank; base=h2

\section{第五章}

马丁使一名强盗归信。\par

从那时起，\Person{马丁}退出军役，热切寻求与\Place{皮克塔瓦城}的主教\Person{希拉略}交往；希拉略在属神之事上的信德在当时备受赞誉且普受尊敬。马丁与他同住了一段时间。这位希拉略授予他执事之职后，曾试图更紧密地将他与自己联结，并引导他参与神圣礼典来约束他。但马丁不断拒绝，高喊自己不配；希拉略作为深具洞察力的人，意识到只能以此方式约束他：若把一种在履行时似乎会对他造成损伤的职务加于他。于是他任命马丁为驱魔者。马丁没有拒绝这项任命，因为担心自己会因轻视这卑微的职务而被看作高傲。不久之后，他在梦中得到警告，应当探访自己的故乡，特别要关心其父母的宗教利益，他们仍陷于异教之中。他遵照圣希拉略明确的意愿出发，希拉略以许多祈祷和泪水恳求他届时返回。据报告，马丁在忧郁的心情中踏上旅程，并召唤弟兄们作证，称有许多苦难在他面前。结果完全证实了这一预言。因为，首先，他在\Place{阿尔卑斯山}间跟随了一些曲折小径，落入强盗之手。其中一个强盗举起斧子悬在马丁头上时，另一个用右手挡住了下落的打击；尽管如此，马丁的双手被绑在背后，被交给其中一个强盗看守并剥去衣物。那强盗将他带到远离其他人的隐秘处，开始问他是什么人。马丁回答说他是基督徒。强盗接着问他是否害怕。这时马丁极其勇敢地回答说，他从未感到如此安全，因为他知道主的仁慈在考验中必将特别与他同在。他补充说，他更怜悯那个看守他的人，因为过着抢劫的生活，显出自己不配得到基督的仁慈。然后他开始讲论福音真理，向强盗宣讲天主的圣言。我何必拖延陈述结果？那强盗相信了；在表达对马丁的敬意后，他将他带回路上，恳求他为主祈祷。那强盗后来被人看见过着虔诚的生活；事实上，我上面叙述的故事是基于他本人提供的记述。\par

% structure: p0019 source=h2 target=section reason=relative-heading-rank; base=h2

\section{第六章}

魔鬼投身在\Person{马丁}的路途。\par

\Person{马丁}从那里继续前行，经过\Place{米兰}后，魔鬼化作人形在路上遇见他。魔鬼先问他去往何处。马丁回答说他愿意去主所召唤的任何地方，魔鬼对他说：\Quote{无论你去哪里，或尝试什么，魔鬼都将抵抗你。}于是马丁用先知的话回答说：\Quote{主是我的助佑；我不惧怕人能对我做什么。}随即，他的仇敌立刻从他眼前消失；这样，正如他心中所愿，他使自己的母亲脱离了异教的错误，虽然他的父亲仍执著于其罪恶。但他以榜样拯救了许多人。\par

此后，当亚略异端蔓延全世界，在伊利里亚尤其猖獗，而他几乎孤身一人与司铎们的背信弃义奋力抗争，并遭受了许多惩罚（因为他被公开鞭打，最终被迫离开城市）时，他再次前往意大利，发现高卢两部分的教会因圣依拉略的离去而陷入混乱——圣依拉略被异端者的暴力驱逐流放——于是他在米兰为自己建立了一座隐修院。在那里，亚略人的首领奥森修也残酷地迫害他，在多次凌辱他之后，暴力将他逐出城市。因此，他认为有必要顺应形势，便与一位品性卓绝的司铎作为同伴，退隐到加里纳里亚岛。他在那里以植物根茎为生，并曾食用藜芦，据说这是一种有毒的草。但当感到体内毒力增强、死亡临近时，他藉着祈祷避开了迫在眉睫的危险，即刻所有疼痛都消失了。不久后，他得知因国王悔罪，圣依拉略获准返回，便设法在罗马与他相见，为此他启程前往该城。\par

% structure: p0023 source=h2 target=section reason=relative-heading-rank; base=h2

\section{第七章}

玛尔定使一位望教者复活\par

依拉略已经离开，玛尔定便追随他的足迹；依拉略极其欢欣地接待了他，玛尔定便在离城不远处为自己建立了一座隐修院。当时，有一位望教者加入了他，希望学习这位最圣洁之人的教义和生活方式。但仅过了几天，这位望教者患上虚弱，开始发高烧。恰巧玛尔定当时离家，外出三天后返回时，发现望教者已经去世；死亡来得如此突然，以至于他未及领受圣洗就离开了人世。尸体被停放在外，哀悼的弟兄们正以最后的悲伤仪式表达敬意，这时玛尔定流着泪、悲叹着赶到他们面前。然后，他仿佛以全部心力抓住了圣神，命令其他人离开停尸的小屋；他闩上门，全身伏卧在已故弟兄的遗体上。他恳切祈祷了一段时间，藉着天主的圣神察觉到力量临在，便稍作起身，凝视着死者的面容，毫不怀疑地等待着祈祷和主仁慈的结果。过了不到两个时辰，他看见死者的所有肢体开始微微活动，眼睛张开，似乎要恢复视力。于是他大声向主感恩，呼喊声充满了小屋。听到响声，站在门口的人立刻冲进来。他们目睹了奇妙的景象：他们看见先前已死的人如今活着。这样，他复活了，随即领受了圣洗，之后又活了多年；他是第一个向我们展示玛尔定德能事迹并为之作证的人。此人常讲述，当他灵魂离体时，被带到审判者的法庭前，被分配到阴暗的境域和粗俗的人群中，并受到严厉的判决。然后，他补充说，有两位天使提示审判者，他是玛尔定为其祈祷的人；因此，他被命令由那两位天使领回，交给玛尔定，恢复先前的生活。从那时起，这位圣者的名声大噪，被所有人视为圣人，也被认为大有能力、真正有德。\par

% structure: p0026 source=h2 target=section reason=relative-heading-rank; base=h2

\section{第八章}

玛尔定使一位自缢者复活\par

此后不久，当玛尔定经过一位名叫卢皮奇努斯者的庄园时——此人在世俗眼中备受尊敬——他听到悲痛的哭喊声。他忧心忡忡地走近人群，询问哭泣的原因，得知这家的一名奴隶上吊自尽了。玛尔定听闻，进入停尸的小屋，屏退众人，俯伏在尸体上，祈祷了片刻。不久，死者脸上显现生命迹象，垂下的眼睛凝视着玛尔定的面容，他苏醒了；他轻轻尝试起身，抓住了圣者的右手，藉此站了起来。就这样，在众人注视下，他与玛尔定一同走到房屋的门廊。\par

% structure: p0029 source=h2 target=section reason=relative-heading-rank; base=h2

\section{第九章}

玛尔定所受的高度敬仰\par

几乎在同一时期，马丁被召担任图尔教会的主教职务；但当他难以从隐修院中被请出来时，一位名叫鲁里基乌斯的市民，假称其妻患病，跪在他膝前，说服他外出。众多市民事先已在他途经的道路上布设了岗哨，他于是像被押送一般被护送至城中。不仅有来自该城，还有来自邻近城市的难以计数的人群，奇妙地聚集而来投票。众人只有一个愿望，同样的祈祷，同样的坚定信念：马丁最配得上主教之职，教会因这样的司铎而蒙福。然而，有少数人，其中包括几位被召来任命首席司铎的主教，竟不虔敬地加以抵制，妄称马丁其人不值一提，不配担任主教，相貌可鄙，衣着寒酸，头发令人作呕。而见识更正确的人们则嘲笑他们的这种愚妄，因为这些反对者越是诽谤，反而越彰显出那人的卓越品格。事实上，除了民众所愿成就的、承行天主旨意的事，他们什么也不能做。在与会的主教中，据说一位名叫德芬索的主教特别反对；因此有人注意到，他在先知诵读中当场受到了严厉斥责。因为那天负责公开诵读的读经者恰被民众挡住未能出现，官员们陷入混乱，等候那位始终未到的人，这时一位旁立者拿起《圣咏集》，抓起了映入眼帘的第一节。那篇圣咏这样写道：\BibleQuote{从婴孩和吃奶者的口中，你建立了赞颂，为使你的仇敌沉默，为使仇敌和复仇者止息。}这些话语诵读出来后，民众发出欢呼，反对派则狼狈不堪。人们相信这篇圣咏是出于上主的安排，为使德芬索听到与他作为相称的见证，因为对马丁而言，上主的赞颂从婴孩和吃奶者的口中得以完成，同时仇敌也被指明并消灭了。\par

% structure: p0032 source=h2 target=section reason=relative-heading-rank; base=h2

\section{第十章}

马丁任图尔主教。\par

如今他担负起主教之职后，我实在无法充分描述马丁如何在履行职责中显扬自己。因为他以最大的恒心，始终如一，与先前一样。他的内心同样谦卑，衣着同样朴素。他兼具威严与和蔼，恰当地保持了主教的地位，却又不放弃隐修士的目标与美德。因此，他有一段时间使用了与教堂相连的小屋；但后来，当他感到无法忍受众多来访者带来的干扰时，便在城外约两里处为自己建立了一座隐修院。这地方如此隐秘僻静，使他能享受独修者的宁静。因为一边是高山陡峭的岩石环绕，卢瓦尔河则通过一段短距离的港湾将平原的其余部分封闭起来；这地方只有一条非常狭窄的通道可以接近。他在那里拥有一间木造的小屋。许多弟兄也以同样的方式为自己建造了退隐之所，但多数人是从悬垂的山岩中开凿洞穴而成。共有八十位门徒，以圣德导师为榜样接受训练。那里没有人拥有称为自己之物；一切皆为公有。不允许买卖任何东西，如同大多数隐修士的惯例。那里不从事任何手艺，除了抄写的工作，而且这工作只分配给较年轻的弟兄，年长的则专心祈祷。很少有人离开自己的小屋，除非是聚集到祈祷处。他们都在禁食时刻过后才一起进餐。没有人饮酒，除非因病被迫如此。多数人身着骆驼毛的衣服。\BibleRef{玛 3:4}任何近乎柔软的服装在那里都被视为有罪，这更显得非凡，因为其中许多人是被认为是贵族出身的人。这些人虽受过迥然不同的教养，却迫使自己降到如此谦卑和忍耐的程度，我们后来看到其中不少人成为了主教。因为哪一个城市或教会不渴望自己的司铎出自马丁的隐修院呢？\par

% structure: p0035 source=h2 target=section reason=relative-heading-rank; base=h2

\section{第十一章}

马丁摧毁奉祀强盗的祭台。\par

但我现在要叙述\Person{马丁}作为主教所显示的其他德行。离镇不远，有一处非常靠近隐修院的地方，曾被一种虚假的人的意见所祝圣，以为有些殉道者合葬在那里。因为人们也相信，以前的主教曾在那里安放了一座祭台。但\Person{马丁}不倾向于轻信不确定之事，常常向那些年长的人——无论是司铎还是神职人员——询问，是否应将那位殉道者的名字或他受难的时间向他说明。他说他这样做，是因为在这些事上他心中有很大的疑虑，因为从古代流传下来关于这些事的可靠传统并不存在。因此，他一时之间远离了那地方，绝不愿意削弱人们对它的宗教敬仰，因为他尚不确定；但同时他也没有将自己的权威借给民众的意见，以免单纯的迷信获得更稳固的基础。有一天，他带着几位弟兄作为同伴，去了那地方。他站在坟墓正上方，向主祈祷，求主启示那人是谁，他的品性或功德如何。然后转向左边，他看见一个面貌卑劣、形态可怖的阴魂站在很近的地方。\Person{马丁}命令他说出自己的名字和品性。于是，他说出了自己的名字，并承认了自己的罪过。他说他曾是一个强盗，因罪行被斩首；他只是因民众的错误而受到尊敬；他与殉道者毫无共同之处，因为光荣是殉道者的份，而惩罚却从他身上索取代价。旁边的人奇妙地听到了说话者的声音，却看不见任何人。然后\Person{马丁}把他所见的讲了出来，并命令把那里的祭台移走，从而将民众从那种迷信的错误中解救出来。\par

% structure: p0038 source=h2 target=section reason=relative-heading-rank; base=h2

\section{第十二章}

\Person{马丁}令抬尸者停下。\par

此后不久，有一次\Person{马丁}外出旅行，遇到一个外教人的尸体正被以迷信的葬礼送往坟墓。他从远处看见迎面而来的队伍，不知道发生了什么事，便稍停片刻。因为他与那队伍相距约半英里，很难看清他目睹的景象究竟是什么。然而，由于他看到那是一个乡民聚集，而盖在尸体上的亚麻布被风吹动，他便以为正在举行某种亵渎的祭献仪式。他之所以这样想，是因为高卢乡民在他们的可怜愚昧中，习惯用白色覆盖物裹着魔鬼的偶像在田野间游行。因此，\Person{马丁}举起十字圣号对着他们，命令那队伍不要移动，并放下所抬的重物。于是，人们可以看到那些可怜的人起初变得像石头一样僵硬；接着他们竭尽全力想要前进，却一步也迈不出，便以最可笑的方式旋转起来，直到再也无法承受重量，只得放下尸体。他们惊愕不已，茫然地互相凝视，不知道发生了什么事，沉默地陷在思想中。但当这位圣德之人发现他们只是一群举行葬礼的农民，而非向诸神祭献时，他便再次举起手，准许他们离开并抬起尸体。这样，他既按自己的旨意迫使它们停下，又在自己认为合适时允许他们离去。\par

% structure: p0041 source=h2 target=section reason=relative-heading-rank; base=h2

\section{第十三章}

\Person{马丁}从倒下的松树中脱险。\par

又一次，当他在某个村庄拆毁了一座非常古老的神庙，并着手砍伐一棵紧邻神庙的松树时，那地方的大司祭和一群其他外教人开始反对他。这些人，虽然由于主的影响，在神庙被拆毁时一直保持安静，却不能容忍这棵树被砍伐。玛尔定仔细地教导他们，树干中没有任何神圣之物，并敦促他们宁可尊敬他自己所事奉的天主。他补充说，砍伐这棵树有道义上的必要，因为它已经被奉献给了一个邪魔。\newline{}那时，他们中一个比别人更胆大的人说：\Quote{如果你对你所声称敬拜的天主有任何信赖，我们自己去砍这棵树，而你要在它倒下时承受它；因为若如你所说，你的主与你同在，你必不受任何伤害。} 于是玛尔定勇敢地依靠主，承诺他会照所请求的去做。于是，所有那一群外教人都同意了所提的条件；因为他们认为失去那棵树是小事，只要能让他们的宗教的敌人被压在树下。因此，因为那棵松树向一边倾斜，所以毫无疑问它会朝哪边倒下，玛尔定被捆绑后，按照这些外教人的决定，被放在那个没有人怀疑树会倒下的地方。于是他们开始非常高兴地砍自己的树，而在远处有大量好奇的围观者。现在那棵松树开始摇晃，并以其倒下的方式威胁自身的毁灭。远处的隐修士们脸色苍白，被越来越近的危险吓坏了，失去了所有的希望和信心，只等待着玛尔定的死亡。但他依靠主，勇敢地等待着，当那倒下的松树发出临终的轰鸣声，正在倒下，正要砸向他时，他仅仅举起手对着它，向它展示救恩的标记。然后，确实，像陀螺一样（人们可能会以为它被推了回去），它向相反的一侧转去，几乎砸碎了那些自以为站在安全地方的人。然后，确实，一声呼喊上达于天，外教人因奇迹而惊异，隐修士们喜极而泣；基督的名被众人共同赞颂。众所周知的结果是，那天救恩临到了那地区。因为在那海量的外教人中，几乎没有一个不表示希望接受覆手，并放弃他不敬的错谬，宣认信仰主耶稣。确实，在玛尔定的时代之前，那些地区很少有人，甚至几乎没有人接受基督的名；但因他的德行和榜样，那名字如此盛行，以至于现在那里没有一个地方不是挤满了教堂或隐修院。因为无论他在哪里摧毁外教神庙，他都会立即在那里建造教堂或隐修院。\par

% structure: p0044 source=h2 target=section reason=relative-heading-rank; base=h2

\section{第十四章}

玛尔定摧毁外教神庙和祭坛。\par

大约在同一时期，他在其他类似的事务中也毫不逊色。因为，他在某个村庄点燃了一座非常古老且著名的神庙，火圈被风带到了紧邻甚至连接着神庙的一座房屋上。玛尔定察觉后，迅速爬上屋顶，站在前进的火焰之前。那时，确实可以看到火焰以奇妙的方式逆风被推回，以至于两种元素之间似乎有一种冲突在相互搏斗。就这样，藉玛尔定的影响，火只在它被命令的地方起作用。\newline{}但在一个名为莱普罗苏姆的村庄，当他同样想推倒一座因对其神圣性的迷信想法而获得巨大财富的神庙时，一群外教人如此强烈地抵抗他，以至于他受了伤而被迫后退。于是，他撤到附近的一个地方，在那里三天，穿着苦衣，蒙灰，禁食并祈祷，恳求主，既然他无法靠人力推倒那座神庙，愿以神能摧毁它。然后两位天使，像天兵一样手持长矛盾牌，突然出现在他面前，说他们是主派来驱散乡民群众，并为玛尔定提供保护，以免在神庙被摧毁时有人抵抗。他们告诉他回去完成他已经开始的蒙福工作。于是玛尔定回到村庄；当外教人群完全安静地看着他将外教神庙夷为平地时，他也将所有祭坛和偶像化为尘土。看到这一幕，那些乡民们发现他们被天主旨意的干预如此惊骇恐惧，以至于不敢再与主教为敌，几乎所有人都信了主耶稣。然后他们开始公开呼喊，承认玛尔定的天主当受敬拜，而那些不能帮助他们的偶像应当被轻视。\par

% structure: p0047 source=h2 target=section reason=relative-heading-rank; base=h2

\section{第十五章}

玛尔定向刺客伸出脖颈。\par

我还要叙述在埃杜伊人村落发生的事。玛尔定在那里推翻一座神庙时，一群乡村异教徒狂怒地向他冲来。其中一人胆大，拔剑攻击玛尔定；玛尔定却掀开披肩，将赤裸的颈项伸给刺客。那异教徒毫不迟疑地要砍下，但在他举起右臂的瞬间，便仰面倒地，被天主的畏惧所压倒，乞求宽恕。类似的事也发生在玛尔定身上：当他摧毁一些偶像时，有一个人决心用刀刺伤他，但就在挥刀之际，武器竟从他手中被打落而消失了。此外，当异教徒恳求他不要推翻他们的神庙时，他常以神圣的言论安抚并化解他们的心，以至真理之光启示他们后，他们亲自推倒了自己的神庙。\par

% structure: p0050 source=h2 target=section reason=relative-heading-rank; base=h2

\section{第十六章}

圣玛尔定所行的治愈\par

此外，玛尔定拥有极大的治愈恩赐，几乎没有任何病人来求助而不立即恢复健康。以下事例清楚表明这一点。在\Place{特里尔}，有一个女孩因严重的瘫痪而完全卧床已久，身体任何部分都不能使用，仿佛已死，只剩一丝气息。她悲痛的亲属站在旁边，只等死亡来临。突然有消息说玛尔定来到了该城。女孩的父亲得知后，跑去为那几乎无生气的孩子恳求。其时玛尔定已进入教堂。在那里，当着众人和许多主教的面，老人悲叹一声，抱住圣人的双膝说：\Quote{我的女儿因一种可悲的病痛快要死了；更可怕的是，她现时仅灵魂活着，肉身早已未死先亡。我恳求您去她那里，为她祝福；因我相信通过您，她将康复。}玛尔定被这话困扰，不知所措，退缩着说此事非他所能；老人判断错了；他不配作主借以显示其大能的器皿。父亲流泪坚持更恳切地恳求，请玛尔定去看看垂死的女孩。最后，在周围主教的敦促下，他去了女孩的家。一大群人在门口等着，要看主的仆人会做什么。首先，他照例使用这类事务中熟悉的武器——俯伏在地祈祷。然后凝视患病的女孩，要求给她油。他接过油并祝福后，将具有大能的圣油倒进女孩口中，她的声音立即恢复。接着，通过与她的接触，她的四肢逐渐恢复生机，直至最后，在众人面前，她迈着坚定的步伐站了起来。\par

% structure: p0053 source=h2 target=section reason=relative-heading-rank; base=h2

\section{第十七章}

玛尔定驱逐多个魔鬼\par

同时，一位曾任总督的\Person{泰特拉狄乌斯}的仆人被魔鬼附身，遭受极其悲惨的折磨。玛尔定应邀为之覆手，便命人将那仆人带来；但恶灵无论如何不肯从仆人所住的房间出来，它张着可怕的利齿，令试图靠近的人胆寒。于是\Person{泰特拉狄乌斯}俯伏在圣者脚下，恳求他亲自去那关押着附魔者的屋子。但玛尔定说不能去一个未皈化异教徒的家，因为\Person{泰特拉狄乌斯}当时仍陷于异教谬误中。\Person{泰特拉狄乌斯}便保证，若魔鬼被逐出那少年，他就做基督徒。玛尔定于是覆手少年，驱走了恶灵。\Person{泰特拉狄乌斯}见状，信了主耶稣，立即成为慕道者，不久后领受了圣洗；他始终对玛尔定怀有特别的深情，视其为救恩的缔造者。\par

大约同时，他进入同城某户人家，刚跨过门槛便停住，说他察觉院内有一个可怕的魔鬼。玛尔定命它离去，它却抓住住在内室的一个家人；那可怜人立即开始用牙齿疯狂地撕咬遇到的人。屋内大乱，家人惊慌，在场的人四散奔逃。玛尔定冲到那狂乱者面前，先命令他站住。但他仍咬牙切齿，张开大口要咬；玛尔定将手指伸进他嘴里说：\Quote{你若有什么能力，就把这些吃掉。}但那人像是烧红的铁进了颚中，远远地缩回牙齿，不敢触碰圣者的手指；而魔鬼被迫以刑罚和折磨离开那附身者，但因无法从口逃出，便从腹中泻出，留下可憎的痕迹。\par

% structure: p0057 source=h2 target=section reason=relative-heading-rank; base=h2

\section{第十八章}

玛尔定行各种奇迹\par

在此期间，一个突然的报告使全城因蛮族的动向与入侵而陷入恐慌。\Person{马尔定}命人将一个附魔者带到他面前，并命令他宣布这消息是否属实。于是他承认有十六个恶魔在人群中散布了这个报告，企图借此引起的恐惧迫使\Person{马尔定}逃离该城；而事实上，蛮族根本没有丝毫入侵的意图。当污灵在教会中这样承认时，城市从当时感受到的恐惧与骚动中获得了释放。\par

又一次，在\Place{巴黎}，当\Person{马尔定}在众多人群的簇拥下进入城门时，他亲吻了一个面色凄惨的癞病人，而众人看到他的举动无不战栗；\Person{马尔定}祝福了他，结果他立刻被治愈，摆脱了所有痛苦。次日，那人皮肤康健地出现在教会中，感谢他所恢复的健康。这一事实也不应默然不语：从\Person{马尔定}的衣服上扯下的线，或是从他穿的苦衣上取下的布条，常在病人身上显奇迹。因为无论是缠在手指上，还是挂在脖子上，它们常常驱散患者的疾病。\par

% structure: p0061 source=h2 target=section reason=relative-heading-rank; base=h2

\section{第十九章}

\Person{马尔定}的一封信治愈疾病，及其他奇迹。\par

此外，\Person{阿尔博里乌斯}，一位前任总督，且是一位极其神圣而忠信的人，他的女儿正因四日疟的烧热而痛苦挣扎时，他在那高热发作之际，将一封刚好送到他手中的\Person{马尔定}的书信塞进了女孩的衣襟，热病立刻消散。此事对\Person{阿尔博里乌斯}影响如此之大，以至于他立即将女儿奉献给天主，并把她许配于永恒的童贞。然后，他前往\Person{马尔定}那里，把女孩呈献给他，作为他行奇迹能力的明显活证，因为她虽不在场却被他治愈了；他决意不让别人，只让\Person{马尔定}亲自给她穿上象征童贞的衣服，以行奉献礼。\par

同样，\Person{保利努斯}，一个日后要成为时代鲜明榜样的人，他的眼睛开始剧痛，且一层相当厚的皮膜长出来，已经遮住了瞳孔；\Person{马尔定}用画家的刷子触摸了他的眼睛，所有疼痛消失，眼睛恢复了原有的健康。\Person{马尔定}自己也曾不慎从楼上的房间摔下，滚过一段破损不平的楼梯，受了多处伤，躺在自己的小室里濒临死亡，忍受着剧痛；他在夜间看见一位天使显现，为他清洗伤口，并给他的伤处涂抹愈合的药膏。结果，次日早晨他发现自己恢复了健康，仿佛从未受过任何伤害一般。但细述此类事情未免冗长，就让这些例子作为众多事例中的少数代表吧；我们不应在显著的神迹干预事件中有损真理，同时，从众多事例中选择这些，也是为了避免引起读者的厌倦。\par

% structure: p0065 source=h2 target=section reason=relative-heading-rank; base=h2

\section{第二十章}

\Person{马尔定}如何对待\Person{马克西穆}皇帝。\par

在此，在如此伟大的事迹中插入一些较小的事（尽管我们时代的本性如此，所有事物都已陷入腐朽和败坏，但司铎的坚定不向皇室的谄媚屈服，几乎是一种卓越的美德），当许多主教从各地聚集到暴烈的皇帝\Person{马西摩}那里时，他因在内战中获胜而洋洋得意，周围所有人的可耻谄媚普遍可见，而司铎的尊严已堕落屈服于皇室的随从之下，唯独在\Person{玛尔定}身上，宗徒的权威继续彰显。因为即使他必须为某些事向皇帝陈情，他也是命令而非恳求；尽管多次被邀请，他仍避开他的宴席，说他不能与一位剥夺了两位皇帝中一人王国、另一人性命的人同席。最后，当\Person{马西摩}坚称他不是自愿夺取王权，而是以武力保卫帝国必要的要求，这是士兵们按天主的安排强加于他的，并且天主的恩宠似乎并未对他缺失，因为他以看似难以置信的方式赢得了胜利，并补充说他的对手中无一人阵亡于公开战场之外时，\Person{玛尔定}终于被他的理由或恳求说服，参加了皇宴。国王因达成这一目的而异常欣喜。此外，在场的宾客如同节日般被邀请：最高贵最显赫的人物——身为执政官的省长，名叫\Person{埃沃迪乌斯}，是世上最正义的人之一；两位权势极大的廷臣，国王的兄弟和叔父；在这两人之间，\Person{玛尔定}的司铎坐下了；而他自己则坐在紧邻国王的座位上。宴席过半，按惯例，一名侍从向国王呈上酒杯。他命令将杯递给最圣洁的主教，期望并希望从他手中接过杯。但\Person{玛尔定}饮完后，将酒杯递给了自己的司铎，认为除了自己之外无人更配饮下一口，并认为他不应把国王本人或国王身边的人置于司铎之上。皇帝及当时在场的所有人都非常钦佩这一举动，以至于他们虽被轻视，却因此感到愉快。于是整个宫中传言，\Person{玛尔定}在国王的宴席上做出了连最低级法官的宴席上也没有主教敢做的事。\Person{玛尔定}曾很早预言同一个\Person{马西摩}，如果他前往当时他想去的意大利，与皇帝\Person{瓦伦提尼安努斯}交战，他将会在初次攻击中获胜，但不久后就会灭亡。我们已看到这事果然发生。因为，当他初到时，\Person{瓦伦提尼安努斯}不得不逃走，但约一年后恢复力量，\Person{马西摩}在\Place{阿奎莱亚}城内被俘并被杀。\par

% structure: p0068 source=h2 target=section reason=relative-heading-rank; base=h2

\section{第二十一章}

\Person{玛尔定}与天使和魔鬼打交道。\par

众所周知，他经常看见天使，他们轮流与他用固定的话语交谈。至于魔鬼，\Person{玛尔定}认为他是如此可见且始终在他眼目之下，以至于无论魔鬼保持自己的形状，还是变成各种属灵邪恶的形象，\Person{玛尔定}都能认出他无论以何种伪装出现。魔鬼深知自己无法逃脱察觉，因此常常辱骂\Person{玛尔定}，因为无法用诡计欺骗他。有一次，魔鬼手持一只血淋淋的公牛角，大声冲进\Person{玛尔定}的房间，伸出他血红的右手，同时为他所犯的罪行而得意，说：\Quote{\Person{玛尔定}，你的能力在哪里？我刚杀了你的一名随从。}于是\Person{玛尔定}召集弟兄们，告诉他们魔鬼所揭露的事，并命令他们仔细搜查各个房间，以查明谁遭遇了这场灾祸。他们报告说，没有一个隐修士失踪，但一个被他们雇用的农夫曾去森林用车拉柴。听到这事，\Person{玛尔定}指示几个人去迎接他。他们去后，发现那人离隐修院不远，几乎已死。尽管如此，他虽即将断气，仍向弟兄们说明了他受伤和死亡的原因。他说，当他拉紧拴在一起的两头牛上松开的皮带时，其中一头牛挣脱了头，用角伤了他的腹股沟。不久那人便断气了。你们看，主以怎样的判断将这种能力赋予了魔鬼。这是\Person{玛尔定}身上可惊奇的一点，不仅在我特别指出的这一场合，而且在许多类似场合，事实上每当这样的事发生时，他都能预先知道，并将所启示的事告诉弟兄们。\par

% structure: p0071 source=h2 target=section reason=relative-heading-rank; base=h2

\section{第二十二章}

\Person{玛尔定}甚至向魔鬼宣讲悔改。\par

当时，魔鬼虽以千百种恶毒手段试图陷害这位圣人，却常常以可见的形式、但极其多样的形状出现在他面前。因为有时他以朱庇特的模样显现，又常以墨丘利和密涅瓦的模样显现。也常听到辱骂之言，一群恶魔用粗鄙的言辞攻击\Person{玛尔定}。但他知道这一切都是虚假和无根据的，并不受所控之罪的影响。此外，有些弟兄作证说，他们曾听见一个恶魔用辱骂的言辞斥责\Person{玛尔定}，问他为何在那些曾因陷入各种错误而失去圣洗的弟兄们后来悔改时，又接纳了他们。那个恶魔列举了他们每个人的罪行；但他们补充说，\Person{玛尔定}坚定地抵抗魔鬼，回答他说，过去的罪过可以通过过更好的生活来洁净，并且因天主的仁慈，那些放弃恶行的人必得赦免。魔鬼反驳说，像他所指出的那些有罪之人不在赦免之列，主不会怜悯那些曾经堕落的人；据说\Person{玛尔定}大声说了如下的话：\Quote{你这个可怜的家伙，若你停止攻击人类，甚至在这审判之日临近的时刻，只为你的行为悔改，我，怀着对主的真正信心，会向你应许基督的怜悯。}啊！这是对主慈爱的何等圣洁的勇敢！虽然他不能声称有权柄，却显露出内心的情感！既然我们的论述在此涉及魔鬼及其诡计，虽然此事不直接关于\Person{玛尔定}，但讲述所发生的事似乎并不离题；既因\Person{玛尔定}的德行在某种程度上显现在这一事件中，又因这件值得神迹的事，为了提供警示以防将来发生类似情况，理当记录下来。\par

% structure: p0074 source=h2 target=section reason=relative-heading-rank; base=h2

\section{第二十三章}

一个魔鬼欺骗的案例\par

有一个人，名叫\Person{克拉鲁斯}，是一位最高贵的青年，后来成为司铎，如今因幸福离世而被列于圣人之列。他弃绝一切，投奔\Person{玛尔定}，不久便因最崇高的信德和各种美德而卓著。当时，他在离主教的隐修院不远处为自己建了一座住所，许多弟兄与他同住。有一个名叫\Person{阿纳托利乌斯}的青年，以隐修士的身份虚假地表现出谦逊和纯真，来到他那里，与其他弟兄一起共用公物生活了一段时间。久而久之，他开始声称天使常常与他交谈。因无人相信他的话，他便用某些记号说服了许多弟兄。最后，他竟然宣称天使在他与天主之间来往，并要求被视作先知之一。但\Person{克拉鲁斯}始终不能相信他。于是，他开始以天主的愤怒和眼前的苦难威胁\Person{克拉鲁斯}，因为他不相信一位圣人。最后，据说他脱口说出以下的话：\Quote{看，主今夜要从天上赐我一件白衣，我穿上它，将住在你们中间；这对你们将是一个记号，表明我是天主的德能，因为我已接受了天主的衣服。}那时，众人因这宣告而充满期待。到了午夜左右，人们听到急切走动的人声，显然整个隐修院都激动起来。也看见那年轻人所住的小屋被众多灯光照亮；能听到屋内走动者的低语和许多声音的嗡嗡声。随后，当安静下来时，那青年出来，叫了一位名叫\Person{萨巴提乌斯}的弟兄，给他看他所穿的衣服。\Person{萨巴提乌斯}惊异万分，召集其余人，\Person{克拉鲁斯}自己也跑来；取得光亮后，他们仔细查看那衣服。这衣服极其柔软，光彩夺目，闪耀着紫色，但没有人能辨认其性质，或是由何种羊毛制成。然而，当用眼睛或手指更仔细检查时，它似乎只不过是一件衣服。同时，\Person{克拉鲁斯}催促弟兄们恳切祈祷，求主更清楚地指示他们这到底是什么。于是，当夜余下的时间用来唱赞美诗和圣咏。但天破晓时，\Person{克拉鲁斯}想要拉着年轻人的手，带他到\Person{玛尔定}那里，深知他不会受魔鬼任何诡计的欺骗。那时，这可怜的人开始抗拒和拒绝，声称他被禁止出现在\Person{玛尔定}面前。当他们强迫他违心前往时，那衣服从带路者手中消失了。因此，谁能怀疑这再次展示了\Person{玛尔定}的能力，以致魔鬼在\Person{玛尔定}的眼前不能再掩饰或隐藏自己的欺骗呢？\par

% structure: p0077 source=h2 target=section reason=relative-heading-rank; base=h2

\section{第二十四章}

玛尔定受到魔鬼诡计的诱惑\par

又发现，大约在同一时期，\Place{西班牙}有一个年轻人，他藉许多征兆在人民中为自己获得了权威，骄傲到如此地步，竟自称是\Person{厄里亚}。当群众过于轻信了这一点后，他进而说他自己其实就是基督；他甚至在这幻觉中也如此成功，以致一位名叫\Person{路福}的主教竟敬拜他为主。我们后来看到这位主教因此被剥夺了职务。许多弟兄也告知我，同时在东方出现了一个人，自诩是\Person{若望}。由此我们可以推断，既然出现了这类假先知，假基督的来临就在眼前；因为他已经在这些人身上操练邪恶的奥秘。而且，我确实认为这一点不应略过：魔鬼大约在此时期用何等诡计诱惑了\Person{玛尔定}。因为，在某一天，在事先祈祷之后，那恶者本身被一道紫色光芒包围，为了借由所装扮的光辉更容易欺骗人，他身披王袍，头戴镶有宝石和金子的冠冕，鞋子也嵌有金子，面露平静之容，一副全然喜悦的神情，使人不会想到他是魔鬼——他站在正在小室内祈祷的玛尔定身边。圣人起初被他外表所眩，长时间保持深沉缄默。首先打破沉默的是魔鬼，他说：\Quote{玛尔定，你要认出你所看见的是谁。我是基督；正快要降临人间，我先愿向你显现。}玛尔定听了这些话，保持沉默，没有作答，魔鬼便胆敢重复他大胆的声明：\Quote{玛尔定，你既然看见了，为什么迟疑不信？我是基督。}那时，圣神启示玛尔定真理，使他明白那是魔鬼而非天主，他便这样回答说：\Quote{主耶稣并没有预言祂会身穿紫袍、头戴闪亮的冠冕而来。除非祂以受难时的容貌和形态出现，并显明十字架上伤口的标记，否则我不会相信基督已来临。}魔鬼听了这些话，便如烟一般消失，并使小室充满如此令人作呕的气味，留下了他真实身份的明确证据。这件事，正如我刚才所述，确实以我所说的方式发生，我是从玛尔定亲口所述得知的；因此，谁也不要视之为虚构。\par

% structure: p0080 source=h2 target=section reason=relative-heading-rank; base=h2

\section{第二十五章}

\Person{苏比西}与\Person{玛尔定}的交往\par

因为我长久以来听闻他的信德、生活与德行，内心燃起了认识他的渴望，便为了见他而踏上了一段于我而言愉快的旅程。同时，因为我的心灵已被撰写他生平的热望所燃烧，我获取信息部分来自他本人——在我敢于询问他的范围内——部分来自那些与他同住或熟知事实的人。而此时，他接待我时的谦逊与仁慈几乎令人难以置信；他衷心为我欢喜，并在主内喜乐，因为他被我如此看重，以至于我因渴望见他而踏上旅途。我这不配的人！（事实上，我几乎不敢承认）他竟屈尊接纳我与他共融！他甚至亲自端水给我洗手，傍晚还亲自洗我的脚；我也没有足够的勇气抵抗或反对他这样做。事实上，我被他不自觉地散发出的权威所征服，以至于我认为除了顺从安排之外，做任何事都是不合法的。他与我交谈的内容都指向以下几点：应当抛弃世俗的诱惑和负担，以便我们能够自由轻松地跟随主\Person{耶稣}；他还向我极力推荐那位杰出人物\Person{保利诺}的行为，作为当前环境的卓越榜样，我前面已经提到过他。\Person{玛尔定}如此评价他：他变卖所有财产跟随基督，因此表明自己在当代几乎是唯一完全遵守\WorkTitle{福音}训导的人。他强烈坚持此人应成为我们效法的对象，并补充说，当前时代是幸运的，因为拥有如此的信德与德行楷模。因为保利诺富有且拥有许多财产，他按照主的明确旨意变卖所有，施舍给穷人，玛尔定说，他已通过实际证明使那些看似不可能的事成为可能。玛尔定的言谈与对话中蕴含着何等的力量与尊严！他多么活跃、多么实际、在解决与\WorkTitle{圣经}相关的问题时多么敏捷而机敏！因为我知道许多人在这一点上是难以置信的——因为我确实遇到过当我说起这些事时不信我的人——我以\Person{耶稣}和我们作为基督徒的共同希望为证，我从未从除玛尔定以外的任何人口中听到如此的知识与天才展现，或如此美好纯洁的话语范例。但即便如此，所有这些赞美与他所拥有的德行相比是何等微不足道！然而，值得注意的是，在一个不能称为博学的人身上，甚至这一属性[高度智慧]也不缺乏。\par

% structure: p0083 source=h2 target=section reason=relative-heading-rank; base=h2

\section{第二十六章}

言语无法描述\Person{玛尔定}的卓越\par

但现在，我的书必须结束，我的论述必须完结。这并不是因为所有关于\Person{玛尔定}值得说的话已经说尽，而是因为我，就像慵懒的诗人到作品结尾时越来越粗心，被题材的浩瀚所挫败而放弃。因为，虽然他的外在行为在某种程度上可以用言语表达，但我坦承，没有任何语言能够描述他的内心生活和日常操行，以及他始终专注于天上之事的心神。没有人能够充分说明他在禁食和守斋中的恒心和自制，或他在守夜和祈祷中的力量，以及他度过的日日夜夜，没有一刻脱离对天主的服侍，无论是为了贪图安逸还是从事事务。而实际上，他只在身体本性需要时才进食或睡眠。我坦然承认，如果像俗话说的，即使\Person{荷马}本人从阴间上来，他也无法用语言公正地处理这个主题；玛尔定身上的所有卓越都如此超越了语言表达的可能性。从未有一个时辰或片刻过去，其中他不是在实际上从事祈祷；或者，倘若他恰好在做别的事，他也从未让自己的心神脱离祈祷。事实上，就像铁匠的习惯，在工作中敲打自己的铁砧以缓解劳动的疲乏，同样，玛尔定即使看似在做别的事，仍在祈祷。啊，真正有福的人，心中没有诡诈——不判断人，不定人罪，不以恶报恶！他在忍受伤害时表现出如此惊人的忍耐，以至于即使当他身为主教时，也允许自己被最低级的圣职人员无端伤害；他既没有因此将他们撤职，也没有（在他力所能及之内）将他们从自己的友爱中排斥出去。\par

% structure: p0086 source=h2 target=section reason=relative-heading-rank; base=h2

\section{第二十七章}

玛尔定令人惊叹的虔诚\par

从未有人见过他发怒、激动、哀哭或欢笑；他始终保持同一：面容上流露出天上的喜乐，仿佛超越了人性的普通界限。他口中从未有过别的言语，只有基督；心中从未有过别的情感，只有虔诚、和平与温柔的慈悲。他也常常为那些显为辱骂他的人——那些在他过着隐居宁静生活时，以毒蛇般的口舌诽谤他的人——的罪而哭泣。我们确实经历过一些嫉妒他的美德和生命的人，他们憎恨在他身上看到而自己没有、也无法模仿的东西。并且——这邪恶实在值得最深的哀伤和呻吟！——他的诽谤者中，虽然很少，但其中有他的恶意中伤者，据说不是别人，竟是一些主教！然而，这里无需点名，尽管其中许多人仍在对我发泄怒气。我认为足以让，如果他们中有人读到这段记述，并意识到自己被指涉，他或许会有羞愧的恩宠。但另一方面，如果他表现出愤怒，他恰恰就承认自己属于那些被谈论的人，尽管自始至终我或许在想的是另一个人。然而，如果那种人把我与玛尔定这样的人一同列入他们的憎恨，我绝不感到羞愧。我深信这一点，即眼下这部小作品将使所有真正善良的人喜悦。我还要进一步说，如果有人以不信的精神阅读这段叙述，他自己将陷入罪中。我自知，我是因着对事实的信德和对基督的爱而被引导写这些事；并且这样做时，我陈述了众所周知之事，记录了真实之事；并且，我相信，那将要获得天主预备之赏报的，不是读这些事的人，而是相信这些事的人。\par

\SourceRecord{Translated by Alexander Roberts.；Nicene and Post-Nicene Fathers, Second Series；Vol. 11.；Edited by Philip Schaff and Henry Wace.；Buffalo, NY: Christian Literature Publishing Co.,；1894.；<http://www.newadvent.org/fathers/3501.htm>.}
